% VI32-DETAILED-PROBLEM-15
\begin{problem}[VI.32.P15: Local dimension and codimension bridge]\label{prob:vi32-15}
Let \(Y\subseteq X\) be irreducible closed with generic point \(\eta_Y\). Reconcile VI/31 with the present definition by proving
\[
\operatorname{codim}(Y,X)=\dim_{\eta_Y}^{\mathrm{loc}}X.
\]

\textbf{Provenance.} Canonical VI/32 synthesis.
\end{problem}

\ifdefined\FullProblemDossiers
\begin{solution}
By definition in VI/32,
\[
\dim_{\eta_Y}^{\mathrm{loc}}X
=
\dim\mathcal O_{X,\eta_Y}.
\]
Choose an affine neighborhood
\[
U=\Spec A
\]
of \(\eta_Y\), and let \(\mathfrak p\) be the corresponding prime. Then
\[
\mathcal O_{X,\eta_Y}\cong A_{\mathfrak p},
\]
so
\[
\dim_{\eta_Y}^{\mathrm{loc}}X
=
\dim A_{\mathfrak p}
=
\operatorname{ht}(\mathfrak p).
\]

On the other hand, VI/31 proved the affine codimension--height dictionary and its intrinsic global form:
\[
\operatorname{codim}(Y,X)
=
\operatorname{ht}(\mathfrak p).
\]
Therefore
\[
\boxed{
\operatorname{codim}(Y,X)
=
\dim\mathcal O_{X,\eta_Y}
=
\dim_{\eta_Y}^{\mathrm{loc}}X.
}
\]

The distinction is conceptual. In VI/31, the local ring appeared only at the generic point of \(Y\) as a device for
measuring relative position. In VI/32, the same invariant is attached to every point \(x\in X\), not just generic
points of closed irreducible subsets.
\end{solution}
\fi
